% Standalone resolution manuscript generated from solution.md.
% Compile with XeLaTeX; author and review metadata are maintained in Markdown.
\documentclass[10pt,a4paper]{article}
\usepackage[margin=24mm,headheight=15pt]{geometry}
\usepackage{fontspec}
\setmainfont{DejaVuSerif.ttf}[BoldFont=DejaVuSerif-Bold.ttf,ItalicFont=DejaVuSerif-Italic.ttf,BoldItalicFont=DejaVuSerif-BoldItalic.ttf]
\setsansfont{DejaVuSans.ttf}[BoldFont=DejaVuSans-Bold.ttf,ItalicFont=DejaVuSans-Oblique.ttf]
\setmonofont{DejaVuSansMono.ttf}
\usepackage{amsmath,amssymb,unicode-math}
\setmathfont{latinmodern-math.otf}
\usepackage{xcolor,microtype,parskip,needspace,fancyhdr}
\usepackage{bookmark}
\usepackage{xurl}
\hypersetup{colorlinks=true,urlcolor=blue,linkcolor=blue,pdftitle={SP-15:
An infinite fiber of shifted singular-value data},pdfauthor={George
Stepaniants},pdflang={en-GB}}
\urlstyle{same}
\setlength{\emergencystretch}{3em}
\providecommand{\tightlist}{\setlength{\itemsep}{0pt}\setlength{\parskip}{0pt}}
\providecommand{\pandocbounded}[1]{#1}
\setcounter{secnumdepth}{0}
\allowdisplaybreaks[1]
\widowpenalty=10000
\clubpenalty=10000
\pagestyle{fancy}
\fancyhf{}
\fancyhead[L]{\small\sffamily NEGATIVE RESOLUTION / SP-15}
\fancyhead[R]{\small\sffamily 12 September 2026 (UTC)}
\fancyfoot[L]{\parbox[b]{0.9\textwidth}{\fontsize{8}{10}\selectfont Substantial
AI assistance; independent automated review is documented separately.}}
\fancyfoot[R]{\footnotesize\thepage}
\begin{document}
{\raggedright\Large\bfseries SP-15: An infinite fiber of shifted
singular-value data\par}
\vspace{0.4em}
{\normalsize\bfseries George Stepaniants\par}
{\small Department of Computing and Mathematical Sciences, California
Institute of Technology, Pasadena, California, USA\par}
\vspace{0.6em}
This manuscript answers the complete
\href{https://github.com/ajt60gaibb/OpenProblemsInNLA/blob/main/eigenvalues-and-inverse-problems/SP-15/README.md}{SP-15
question} negatively in dimension nine. The original question and
generic finiteness theorem are due to Maxime Fortier Bourque and Thomas
Ransford. The argument below is an existence proof of a continuous
exceptional fiber, using an exact block determinant identity and the
smooth constant-rank theorem.

\textbf{Independent review.} A separate
\href{https://github.com/ajt60gaibb/OpenProblemsInNLA/blob/main/references/stepaniants-sp15-2026-09-12/verification/independent-review-aa01/review.md}{Codex-agent
mathematical audit} returned PASS for the complete negative resolution,
with no mathematical correction. The
\href{https://github.com/ajt60gaibb/OpenProblemsInNLA/blob/main/references/stepaniants-sp15-2026-09-12/README.md}{submission
record} preserves the exact frozen source and signed reports.
Substantial AI assistance is disclosed; this is informal automated
review, not external human peer review or formal verification.

\subsection{Exact target and
conclusion}\label{exact-target-and-conclusion}

SP-15 asks whether for each complex matrix order \(N\) there is a finite
uniform bound on the number of unitary similarity classes sharing all
the singular values of \(A-zI\), for every \(z\in\mathbb C\).

\textbf{Theorem.} There is a nonconstant smooth one-parameter family of
complex \(9\times9\) matrices whose members have the same singular
values for every complex scalar shift, and no two distinct members are
unitarily similar. Each member is nilpotent of index three, with Jordan
type \((3,3,3)\). Consequently no integer \(M_9\) with the property
required in SP-15 exists, and the proposed assertion for every \(N\) is
false.

The construction is an existence argument using a polynomial coefficient
map and the constant-rank theorem. It does not require a numerical
search, and does not assert a specific numerical choice of the fixed
data fiber.

\subsection{1. A three-block family and its shifted
determinant}\label{a-three-block-family-and-its-shifted-determinant}

Let \(r\geq1\), and let \(P,Q\in\mathbb C^{r\times r}\) be positive
definite Hermitian matrices. Put \(X=P^{1/2}\), \(Y=Q^{1/2}\), where
both roots are the positive definite Hermitian roots, and define

\nopagebreak[4]
\[
A(P,Q)=\begin{pmatrix}0&X&0\\0&0&Y\\0&0&0\end{pmatrix}.
\]

For real \(t>0\) and \(z\in\mathbb C\), write \(s=t+|z|^2\) and
\(K=s^2I_r+tP\). Direct block multiplication gives

\nopagebreak[4]
\[
tI_{3r}+(A-zI_{3r})^*(A-zI_{3r})
=\begin{pmatrix}
sI_r&-\overline zX&0\\
-zX^*&sI_r+P&-\overline zY\\
0&-zY^*&sI_r+Y^*Y
\end{pmatrix}.
\]

Eliminating the first block leaves a middle diagonal block

\nopagebreak[4]
\[
sI_r+P-\frac{|z|^2}{s}P
=sI_r+\frac tsP=\frac Ks.
\]

\Needspace{12\baselineskip}
Both \(sI_r\) and \(K\) are positive definite. Taking the next Schur
complement and using \(YY^*=Q\) gives

\nopagebreak[4]
\[
\begin{aligned}
&\det\bigl[tI_{3r}+(A-zI_{3r})^*(A-zI_{3r})\bigr]\\
&\qquad=\det K\,
\det\bigl[sI_r+Y^*(I_r-|z|^2sK^{-1})Y\bigr]\\
&\qquad=\det K\,
\det\bigl[sI_r+(I_r-|z|^2sK^{-1})Q\bigr]\\
&\qquad=\det\bigl[sK+(K-|z|^2sI_r)Q\bigr]\\
&\qquad=\det\bigl[s^3I_r+st(P+Q)+tPQ\bigr].
\end{aligned}
\]

The second equality uses \(\det(sI+BC)=\det(sI+CB)\); the third
multiplies the matrix inside the determinant on the left by \(K\). The
final equality uses \(s-|z|^2=t\). No commutativity of \(P\) and \(Q\)
is required.

Define the two-variable polynomial

\nopagebreak[4]
\[
F_{P,Q}(u,s)=\det\bigl[uI_r+s(P+Q)+PQ\bigr].
\]

The determinant identity becomes

\nopagebreak[4]
\[
\det\bigl[tI_{3r}+(A(P,Q)-zI_{3r})^*(A(P,Q)-zI_{3r})\bigr]
=t^r F_{P,Q}(s^3/t,s),\qquad t>0.
\]

Therefore equality of \(F_{P,Q}\) for two pairs implies equality of the
Gram characteristic polynomials for every \(z\): for fixed \(z\), the
displayed determinant polynomials agree for all \(t>0\), hence
identically. Their roots are the negatives of the squared singular
values, with multiplicities. Thus equal \(F\) implies equality of every
shifted singular value, exactly as required by the definition in SP-15.

\subsection{\texorpdfstring{2. The data consists of at most nine real
numbers when
\(r=3\)}{2. The data consists of at most nine real numbers when r=3}}\label{the-data-consists-of-at-most-nine-real-numbers-when-r3}

The total degree of \(F_{P,Q}(u,s)\) is at most \(r\), and its \(u^r\)
coefficient is identically one. All its coefficients are real. To see
the latter point, for real \(u,s\) write

\nopagebreak[4]
\[
uI_r+s(P+Q)+PQ=(P+sI_r)(Q+sI_r)+(u-s^2)I_r.
\]

Taking the conjugate transpose reverses the two Hermitian factors. The
identity

\nopagebreak[4]
\[
\det(BC+cI_r)=\det(CB+cI_r)
\]

shows that the determinant equals its complex conjugate. A polynomial
which is real on all of \(\mathbb R^2\) has real coefficients.

When \(r=3\), the monomials of total degree at most three are

\nopagebreak[4]
\[
1,\ u,\ s,\ u^2,\ us,\ s^2,\ u^3,\ u^2s,\ us^2,\ s^3.
\]

There are ten monomials and the coefficient of \(u^3\) is fixed. Thus
the other nine real coefficients give a map into \(\mathbb R^9\) which
determines all shifted singular values of \(A(P,Q)\).

\Needspace{22\baselineskip}
\subsection{3. Ten parameters with no repeated unitary
class}\label{ten-parameters-with-no-repeated-unitary-class}

Take \(r=3\) and restrict to

\nopagebreak[4]
\[
P=\operatorname{diag}(p_1,p_2,p_3),\qquad 0<p_1<p_2<p_3,
\]

and

\nopagebreak[4]
\[
Q=\begin{pmatrix}
q_1&a&b\\
a&q_2&c+id\\
b&c-id&q_3
\end{pmatrix}>0,\qquad a>0,\quad b>0,
\]

where all ten parameters

\nopagebreak[4]
\[
(p_1,p_2,p_3,q_1,q_2,q_3,a,b,c,d)
\]

are real. Their allowed set \(\Omega\) is a nonempty open subset of
\(\mathbb R^{10}\). Strict positivity and all strict scalar inequalities
are open conditions. For example, \(P=\operatorname{diag}(1,2,3)\) and
\(Q\) with diagonal entries \(4\) and all off-diagonal entries \(1\)
belong to this set.

We claim that distinct points of \(\Omega\) give matrices \(A(P,Q)\) in
distinct unitary similarity classes. Indeed, since \(X\) and \(Y\) are
invertible,

\nopagebreak[4]
\[
\ker A=\mathbb C^3\oplus0\oplus0,
\qquad
\ker A^2=\mathbb C^3\oplus\mathbb C^3\oplus0.
\]

If \(A(P',Q')=U^*A(P,Q)U\) with \(U\) unitary, then \(U\) preserves both
of these coordinate subspaces. A unitary preserving this nested flag
preserves its three orthogonal summands, so
\(U=\operatorname{diag}(U_1,U_2,U_3)\). The two nonzero blocks then give

\nopagebreak[4]
\[
X'=U_1^*XU_2,\qquad Y'=U_2^*YU_3.
\]

Consequently

\nopagebreak[4]
\[
P'=U_2^*PU_2,\qquad Q'=U_2^*QU_2.
\]

The strict ordering of the diagonal entries of \(P\) and \(P'\) implies
\(P'=P\) and \(p_j'=p_j\) for each \(j\). Their spectrum is simple, so
\(U_2\) is diagonal unitary, say
\(U_2=\operatorname{diag}(e^{i\theta_1},e^{i\theta_2},e^{i\theta_3})\).
The \((1,2)\) and \((1,3)\) entries of \(Q\) and \(Q'\) are positive
real numbers. The identities

\nopagebreak[4]
\[
a'=e^{i(\theta_2-\theta_1)}a,\qquad
b'=e^{i(\theta_3-\theta_1)}b
\]

force both phases to equal one. Hence \(U_2\) is scalar and \(Q'=Q\).
This proves the claim, including all possible unitary intertwiners of
the full \(9\times9\) matrices.

\subsection{4. A positive-dimensional
fiber}\label{a-positive-dimensional-fiber}

Let \(\Phi:\Omega\to\mathbb R^9\) assign the nine non-leading
coefficients of \(F_{P,Q}\). Every entry of \(P\) and \(Q\) is linear in
the ten real parameters, so \(\Phi\) is a real polynomial map, in
particular smooth.

Let \(k\) be the maximum rank of \(D\Phi\) attained on \(\Omega\). This
maximum exists because the nonempty set of attained ranks is a subset of
\(\{0,1,\ldots,9\}\). Choose \(x_0\in\Omega\) at which the rank equals
\(k\). If \(k>0\), a nonzero \(k\times k\) minor stays nonzero on a
neighborhood of \(x_0\); the rank is at least \(k\) there, and is at
most \(k\) by maximality. If \(k=0\), the rank is already identically
zero on \(\Omega\). Thus in either case the derivative has constant rank
\(k\) on a neighborhood of \(x_0\).

By the constant-rank theorem, after smooth changes of coordinates on
this neighborhood and on a neighborhood of \(\Phi(x_0)\), the map has
the form

\nopagebreak[4]
\[
(x_1,\ldots,x_{10})\longmapsto(x_1,\ldots,x_k,0,\ldots,0).
\]

In particular its fiber through \(x_0\) contains a smooth submanifold of
dimension \(10-k\geq1\), and therefore contains an injective smooth
curve parametrized by a nonempty open interval. All parameters on this
curve lie in \(\Omega\), have the same polynomial \(F\), and correspond
to distinct unitary classes by Section 3. Section 1 supplies equality of
every singular value of every complex shift.

The positive definite square-root map is smooth, so applying
\((P,Q)\mapsto A(P,Q)\) to the curve gives the smooth family asserted in
the theorem. Finally \(A^3=0\) and the upper-right block of \(A^2\) is
the invertible matrix \(XY\), so its nilpotency index is three. Its
kernel dimensions are \(3,6,9\), and hence its Jordan type is
\((3,3,3)\).

Taking arbitrarily many distinct points on this curve disproves every
proposed finite bound \(M_9\). This is a counterexample to the full
universal statement, and does not rely on a generic-finiteness
assertion.

\subsection{5. Scope, attribution, and primary
sources}\label{scope-attribution-and-primary-sources}

The question and generic finiteness theorem are due to Maxime Fortier
Bourque and Thomas Ransford. Their result outside a closed measure-zero
exceptional set is compatible with this exceptional nilpotent family.
The proof settles the universal finiteness question negatively; it does
not claim a classification in other dimensions or an explicit numerical
point on the continuous fiber.

The proof is independent of the matrix-theoretic conclusions cited
below. Its only external general theorem is the standard smooth
constant-rank theorem. The statement and proof explicitly supply its
local constant-rank hypothesis.

\begin{itemize}
\tightlist
\item
  M. Fortier Bourque and T. Ransford,
  \href{https://doi.org/10.1112/jlms/jdn085}{\emph{Super-identical
  pseudospectra}}, \emph{Journal of the London Mathematical Society}
  \textbf{79}(2), 511-528 (2009), Theorem 1.4 and Section 6.2. The
  latter asks whether the exceptional set may be empty.
\item
  T. Ransford,
  \href{https://www.impan.pl/shop/en/publication/transaction/download/product/86371}{\emph{Pseudospectra
  and matrix behaviour}}, \emph{Banach Center Publications} \textbf{91},
  327-338 (2010), Theorem 5.4 and its following discussion on printed
  pages 336-337. \href{https://doi.org/10.4064/bc91-0-19}{DOI}.
\item
  G. Armentia, J. M. Gracia and F. E. Velasco,
  \href{https://doi.org/10.1016/j.laa.2011.01.014}{\emph{Identical
  pseudospectra of any geometric multiplicity}}, \emph{Linear Algebra
  and its Applications} \textbf{436}(6), 1683-1688 (2012). Their
  ordinary-similarity theorem is consistent with the fixed Jordan type
  of the family above.
\end{itemize}

The
\href{https://github.com/ajt60gaibb/OpenProblemsInNLA/blob/main/references/stepaniants-sp15-2026-09-12/verification/identity-check.json}{exact
diagnostic} checks five determinant evaluations, including complex
shifts. Those finite checks are supporting diagnostics; the proof of the
full quantified identity and of the continuous fiber is the analytic
argument in Sections 1-4. The
\href{https://github.com/ajt60gaibb/OpenProblemsInNLA/blob/main/references/stepaniants-sp15-2026-09-12/verification/network-check.json}{dated
public eligibility check} and the
\href{https://github.com/ajt60gaibb/OpenProblemsInNLA/blob/main/references/stepaniants-sp15-2026-09-12/verification/eligibility-and-checks.md}{bounded
source-search record} are separate from mathematical verification.
\end{document}
